\section{Experimental Setup}
\label{sec:setup}

\paragraph{Generators and schemes.}
Table~\ref{tab:generators} lists the generators. Both rows are open-weight models, so the sampling distribution is available and the same SynthID reimplementation can be applied to each. The G1 row uses \texttt{facebook/opt-1.3b}, the default model of the MarkLLM SynthID configuration. The G2 row uses \texttt{Qwen/Qwen2.5-1.5B-Instruct}, from a different family. Neither row is hosted: a hosted API exposes no logits, and Google retired text-watermark detection on its API in August 2026. The detector is MarkLLM's SynthID reimplementation, a research implementation with its own keying, and every removal number is relative to it.

\begin{table}[t]
\centering
\small
\caption{Generators and the scheme each row carries. Logits indicates whether the sampling distribution is available to the watermarking implementation.}
\label{tab:generators}
\begin{tabular}{@{}llll@{}}
\toprule
Row & Model & Scheme & Logits \\
\midrule
G1 & \texttt{facebook/opt-1.3b} & MarkLLM SynthID reimplementation & Yes \\
G2 & \texttt{Qwen/Qwen2.5-1.5B-Instruct} & MarkLLM SynthID reimplementation & Yes \\
\bottomrule
\end{tabular}
\end{table}

\paragraph{Attackers.}
Every tactic that calls a model uses a model from a family different from the generator, so that the attacker cannot reintroduce the generator's own mark. Table~\ref{tab:attackers} lists the models and their roles. The rule is applied per generator row: a chain that attacks a G1 document never uses a G1-family model, and the same holds for G2.

One attacker is used per chain, so a chain's result is attributable to the tactic sequence rather than to a model switch inside the sequence. The search and the intensity sweep use DeepSeek Flash. After the search, the top five chains, the best single tactic, and the Unicode control are re-run end to end with Claude Sonnet on the held-out split. That ablation tests whether the ranking transfers across attacker capability, and its cost stays a small fraction of the search. Exact model snapshots and prices are pinned in \texttt{configs/attackers.yaml} and recorded in every run manifest.

\begin{table}[t]
\centering
\small
\caption{Models used by the attack tactics.}
\label{tab:attackers}
\begin{tabular}{@{}>{\raggedright\arraybackslash}p{3.4cm} >{\raggedright\arraybackslash}p{3.2cm} >{\raggedright\arraybackslash}p{7.2cm}@{}}
\toprule
Role & Model & Used by \\
\midrule
Paraphrase and style     & DeepSeek Flash & Instruction paraphrase, discourse paraphrase, style, document rewrite \\
Translation              & DeepSeek Flash & All back-translation tactics \\
Masked infilling         & \texttt{roberta-large} & Masked-LM infilling \\
Synonym candidates       & DeepSeek Flash & Synonym substitution \\
Regeneration             & DeepSeek Flash & Non-origin regeneration \\
Second-attacker ablation & Claude Sonnet & Top five chains, best single, control \\
\bottomrule
\end{tabular}
\end{table}

\paragraph{Data.}
Table~\ref{tab:data} lists the prompt sets. Three domains cover distinct register and length profiles, and the MarkMyWords tasks are included so that at least part of the evaluation can be compared with published tamper-resistance numbers. Each domain contributes 100 prompts. With four domains and three keys, each generator row contributes 1{,}200 documents, and the two rows contribute 2{,}400 before exclusions.

\begin{table}[t]
\centering
\small
\caption{Prompt sets.}
\label{tab:data}
\begin{tabular}{@{}llll@{}}
\toprule
Domain & Source & Prompts & Target length \\
\midrule
News writing        & \ph{source} & 100 & \ph{tokens} \\
Long-form question answering & \ph{source} & 100 & \ph{tokens} \\
Creative writing    & \ph{source} & 100 & \ph{tokens} \\
MarkMyWords tasks   & \citet{markmywords} & 100 & \ph{tokens} \\
\bottomrule
\end{tabular}
\end{table}

\paragraph{Generation settings.}
Decoding uses temperature $0.7$, top-$p$ $0.9$, and $400$ new tokens for every generator, with the same prompt set across rows. Each prompt is generated once per key. Documents below 150 tokens are discarded, because the detector's power depends on length and short documents would be removed trivially. We report the number discarded.

\paragraph{Keys and repetitions.}
Each document is watermarked under $L = 3$ independent keys drawn from the scheme's keyspace and never reused across runs. Detection is evaluated under the key used at generation. Where the tactic is stochastic, each document passes through the tactic three times with different seeds, and the reported score is the mean. Seeds are fixed and recorded in the configuration files.

\paragraph{Precision and power.}
The $60/40$ split leaves 720 documents per generator row for the search and 480 for the held-out report. At a removal rate near $0.5$, a $95\%$ Wilson interval spans about $\pm 4.5$ percentage points on 480 documents and $\pm 3.6$ on 720; a per-domain cell on the held-out split has 120 documents and spans about $\pm 8.9$ points. Averaging three seeds over a binary outcome puts each document's score on a $1/3$ grid, which sets the per-document resolution; the interval on the mean is set by the document count. The preregistered minimum detectable difference is $0.05$ absolute removal rate. A paired comparison of two chains on the same 480 documents reaches $80\%$ power at that difference when the within-document correlation between the two chains' outcomes is about $0.7$ or higher, and less when it is lower; the pilot run estimates the correlation, and we report the shortfall if the design misses the target. A fallback design with two domains at 60 prompts would leave 144 held-out documents per row; there the same difference is below $50\%$ power even at a within-pair correlation of $0.8$.

\paragraph{Oracle.}
Detection uses the SynthID scheme implementation in MarkLLM \citep{markllm}, configured with the same parameters used at generation. The per-document $p$-value comes from a normal approximation to the unwatermarked null, with variance $1/(12\,n_{\mathrm{eff}})$ over the $n_{\mathrm{eff}}$ unmasked positions; we calibrate it on 100 unwatermarked documents per generator row. This is a same-config oracle: it detects the scheme it was given, under the key family it was given. It does not certify what a vendor detector would report, and we make no such claim. At the corpus operating point, 100 negatives admit one false positive, so that point is coarser than the per-document rate. Section~\ref{sec:limitations} states the consequences.

\paragraph{Compute.}
All runs complete on \ph{hardware description} in \ph{wall-clock time}. The attacker-model calls account for \ph{share} of the cost, and the detector passes for \ph{share}. Appendix~\ref{app:hyper} lists the per-tactic settings, and the configuration files in \texttt{configs/} are the authoritative record.
